\section{扭转}\label{sec:5}
\subsection{抽象运算}\label{sec:5a}
固定 $\Ai$ 范畴 $\A$ 及其对象 $Y$。给定 $\A$ 上的 $\Ai$ 模 $\M$，可以如下定义另一个模 $\C=\mathcal T_Y\M$：
\begin{definition*}

\begin{equation}\label{eq:5.1}
\begin{aligned}
\C(X)&=(\M(Y)\otimes\hom_{\A}(X,Y))[1]\oplus\M(X),\\
\mu^1_{\C}(c_0\otimes b,c_1)
&=\bigl((-1)^{|b|-1}\mu^1_{\M}(c_0)\otimes b+c_0\otimes\mu^1_{\A}(b),\\[-2pt]
&\qquad\mu^1_{\M}(c_1)+\mu^2_{\M}(c_0,b)\bigr),\\
\mu^d_{\C}((c_0\otimes b,c_1),a_{d-1},\ldots,a_1)
&=\bigl(c_0\otimes\mu^d_{\A}(b,a_{d-1},\ldots,a_1),\\[-2pt]
&\qquad\mu^d_{\M}(c_1,a_{d-1},\ldots,a_1)\\[-2pt]
&\qquad+\mu^{d+1}_{\M}(c_0,b,a_{d-1},\ldots,a_1)\bigr)\quad(d\geq2).
\end{aligned}
\end{equation}
\end{definition*}
$\C$ 的 c-含幺性可以用与\eref{3.6}相同的方法证明。

映射 $\M\mapsto\mathcal T_Y\M$ 是 $\Ai$ 函子 $\mathcal T_Y:\Q\to\Q$ 的一部分；这里仍令 $\Q=\modcat(\A)$，称此函子为\emph{沿 $Y$ 的抽象扭转函子}。预模同态 $t\in\hom_{\Q}(\M_0,\M_1)$ 的像为
\[
\begin{aligned}
\tilde t&=\mathcal T_Y^1(t)\in\hom_{\Q}(\mathcal T_Y\M_0,\mathcal T_Y\M_1),\\
\tilde t^1(c_0\otimes b,c_1)&=\bigl((-1)^{|b|-1}t^1(c_0)\otimes b,\ t^1(c_1)+t^2(c_0,b)\bigr),\\
\tilde t^d((c_0\otimes b,c_1),a_{d-1},\ldots,a_1)
&=\bigl(0,\ t^d(c_1,a_{d-1},\ldots,a_1)\\[-2pt]
&\qquad+t^{d+1}(c_0,b,a_{d-1},\ldots,a_1)\bigr)\quad(d\geq2).
\end{aligned}
\]
高阶项 $\mathcal T_Y^d$（$d\geq2$）为零，因此 $\mathcal T_Y$ 实际上是 dg 函子。考察\eref{2.10}可见，它还是严格含幺的。

可以将抽象扭转的构造分成两步。首先，取\secref{3c}意义下的张量积 $\M(Y)\otimes\Y$。它带有\emph{抽象求值同态}
\begin{equation}\label{eq:5.2}
\begin{aligned}
\phi&:\M(Y)\otimes\Y\longrightarrow\M,\\
\phi^d(c_0\otimes b,a_{d-1},\ldots,a_1)&=\mu^{d+1}_{\M}(c_0,b,a_{d-1},\ldots,a_1).
\end{aligned}
\end{equation}
而 $\mathcal T_Y\M=\ACone(\phi)$ 就是它在\eref{3.32}意义下的锥。特别地，有相应模同态 $\iota,\pi$，与 $\phi$ 一起构成
\SourcePage{71}
$H(\Q)$ 中的正合三角
\SourceEquationDiagram{diagram-5-3}{0.30946121}{5.3}
\begin{note}\label{st:5.1}
扭转运算很容易推广到有多个扭转对象的情形。例如，给定两个对象 $Y_+,Y_-$ 和 $\Ai$ 模 $\M$，定义 $\C=\mathcal T_{Y_-,Y_+}\M$ 为
\[
\begin{aligned}
\C(X)&=\bigl(\M(Y_-)\otimes\hom_{\A}(X,Y_-)\\[-2pt]
&\qquad\oplus\M(Y_+)\otimes\hom_{\A}(X,Y_+)\bigr)[1]\oplus\M(X),\\
\mu^1_{\C}(c_{0,-}\otimes b_-,c_{0,+}\otimes b_+,c_1)
&=\bigl((-1)^{|b_-|-1}\mu^1_{\M}(c_{0,-})\otimes b_-+c_{0,-}\otimes\mu^1_{\A}(b_-)\\[-2pt]
&\qquad+(-1)^{|b_+|-1}\mu^1_{\M}(c_{0,+})\otimes b_++c_{0,+}\otimes\mu^1_{\A}(b_+),\\[-2pt]
&\qquad\mu^1_{\M}(c_1)+\mu^2_{\M}(c_{0,-},b_-)+\mu^2_{\M}(c_{0,+},b_+)\bigr),
\end{aligned}
\]
高阶项也以类似方式定义。若 $\hom_{\A}(Y_-,Y_+)=0$，那么这与 $\mathcal T_{Y_-}(\mathcal T_{Y_+}\M)$ 是同一个 $\Ai$ 模。若 $\hom_{\A}(Y_-,Y_+)$ 仅仅无上同调，则存在一个模到另一个模的显然嵌入，其商无上同调，因此两种运算仍给出 $H^0(\Q)$ 中同构的对象。
\end{note}

\subsection{Yoneda 模的扭转}\label{sec:5b}
\begin{definition*}
给定 $\A$ 的两个对象 $Y_0,Y_1$，用 $T_{Y_0}(Y_1)$ 表示任意一个拟表示 $\mathcal T_{Y_0}(\Y_1)$ 的对象，并称其为\emph{沿 $Y_0$ 扭转 $Y_1$} 的结果。
\end{definition*}
若它存在，就带有 $H^0(\A)$ 中的典范态射 $[i]:Y_1\to T_{Y_0}(Y_1)$，后者在 Yoneda 嵌入下成为 $[\iota]$。

现在假设 $H(\A)$ 对有限直和及平移 $S^{\sigma}$（$\sigma\in\Z$）封闭；并假设 $\Hom_{H(\A)}(Y_0,Y_1)$ 为有限维分次向量空间。那么由\stref{笔记}{3.4}，可以构造张量积 $\Hom_{H(\A)}(Y_0,Y_1)\otimes Y_0$，并且它带有到 $Y_1$ 的典范求值态射 $[f]$，即 $[\phi]$ 在 Yoneda 嵌入下的原像。具体地，若选取 $\Hom_{H(\A)}(Y_0,Y_1)$ 的一组基 $[b^i]$，令 $\sigma^i=-|b^i|$，并把张量积看作 $Y_0$ 的若干平移的直和，那么 $[f]$ 在第 $i$ 个直和项上的限制，就是 $[b^i]$ 在\eref{3.5}中的同构
\[
\Hom^{-\sigma^i}_{H(\A)}(Y_0,Y_1)\cong\Hom^0_{H(\A)}(S^{\sigma^i}Y_0,Y_1)
\]
\SourcePage{72}
下的像。若 $T_{Y_0}(Y_1)$ 在 $\A$ 中存在，它位于三角
\SourceEquationDiagram{diagram-5-4}{0.42369185}{5.4}
之中，其中 $[p]$ 在 Yoneda 嵌入下对应于 $[\pi]$。应用\stref{推论}{3.10}可知，\eref{5.4}正合。\znote{64}{三角 (5.4) 是扭转到扭曲复形情形的具体化：$T_{Y_0}Y_1$ 由态射 $f:\hom_{\A}(Y_0,Y_1)\otimes Y_0\to Y_1$ 的锥给出，$[i]$ 与 $[p]$ 分别是锥的标准态射与投影。在 Bondal–Kapranov 的语言中，它正是由对象 $Y_1$（次数 $0$）与 $Y_0$（次数 $-|b^i|+1$，$i$ 取遍一组基）组成的扭曲复形，其微分由参数化的态射 $\hom(Y_0,Y_1)\otimes Y_0\to Y_1$ 给出。[来源：nLab "twisted complex"，https:/\allowbreak /\allowbreak ncatlab.org/\allowbreak nlab/\allowbreak show/\allowbreak twisted+\allowbreak complex（2026-11-15 访问）。]}

最后，考虑 $\A$ 三角化且 c-有限（上同调有限）的情形。
\begin{definition*}
所谓 c-有限，是指对于任意两个对象 $Y_0,Y_1$，空间 $\Hom_{H(\A)}(Y_0,Y_1)$ 都有限维。
\end{definition*}
$\A$ 的三角化性质蕴含它对直和与平移封闭（\stref{笔记}{3.5}）。更重要的是，求值态射 $[f]$ 在 $\A$ 中有映射锥，并位于相应的标准正合三角中。由\stref{推论}{3.16}，这个锥在 Yoneda 嵌入下的像同构于 $\mathcal T_{Y_0}(\Y_1)$。因此，所有对象 $T_{Y_0}(Y_1)$ 都在 $\A$ 中存在。利用与\secref{3d}相同的论证，可以使扭转成为函子性的。也就是说，对每个 $Y\in\Ob\A$，存在\emph{扭转函子} $T_Y:\A\to\A$，它在 $H^0(\operatorname{fun}(\A,\A))$ 中的同构意义下唯一，并且使图
\SourceEquationDiagram{diagram-5-5}{0.18484595}{5.5}
在 $H^0(\operatorname{fun}(\A,\Q))$ 中在同构意义下交换。

一个重要特例是：$\B=\Tw\A$ 是 c-有限 $\Ai$ 范畴上的扭曲复形范畴。回顾\stref{引理}{3.36}，Yoneda 嵌入 $\A\to\Q=\modcat(\A)$ 自然延伸为上同调全忠实函子 $\widetilde{\mathcal I}:\B\to\Q$。固定 $Y_0\in\Ob\A$、$Y_1\in\Ob\B$。由定义，Yoneda 嵌入将 $T_{Y_0}(Y_1)\in\Ob\B$ 映为在 $\modcat(\B)$ 中定义的抽象扭转 $\mathcal T_{Y_0}(\Y_1)$。此外，抽象扭转的构造与限制 $\modcat(\B)\to\Q$ 相容。因此，$\B$ 与 $\Q$ 中的扭转概念相容：
\begin{equation}\label{eq:5.6}
\widetilde{\mathcal I}(T_{Y_0}(Y_1))\cong\mathcal T_{Y_0}(\widetilde{\mathcal I}(Y_1))\quad\text{在 }H^0(\Q)\text{ 中}.
\end{equation}
也可以由\eref{5.4}与\stref{推论}{3.16}推出这一点，但上面的论证给出了一个典范同构。这不仅稍微更精确，还可以使该同构具有\eref{5.5}意义下的函子性。
\SourcePage{73}
\begin{note}\label{st:5.2}
将上述讨论推广到\stref{笔记}{5.1}的情形是直接的。所得对象 $T_{Y_{0,-},Y_{0,+}}(Y_1)$ 位于正合三角
\SourceDiagram{diagram-5-2}{0.55453786}
之中。
\end{note}

\subsection{判定拟表示对象}\label{sec:5c}
给定 $\Ai$ 范畴 $\A$ 的对象 $Y_0,Y_1,Y_2$，若有一个显式判据，能够判断何时可以取 $Y_2$ 为 $\mathcal T_{Y_0}\Y_1$ 的拟表示对象，将会很有用。可以以\stref{引理}{3.1}为基础，因为它指明了定义态射 $\Y_2\to\mathcal T_{Y_0}\Y_1$ 所需的数据；不过，我们更愿意考虑反方向的映射。这一选择以及下面表述中的技术细节，都由预期的几何应用决定（见\sr{sec:17}{第 17 节}）。\znote{65}{之所以考察“从 $\mathcal T_{Y_0}\Y_1$ 到 $\Y_2$ 的映射”而不是反方向，是因为几何应用中真正给出的是 $Y_2$ 与 $T_{Y_0}Y_1$ 拟同构的一个显式代表：清空循环与伪全纯圆盘的计数自然给出一次态射 $c:Y_1\to Y_2$ 以及把 $\hom(Y_0,Y_1)$ 提升一阶的映射 $k$，即数据 $(c,k)$。若反过来写，同样的几何数据就要摊成无限多个方程。[来源：译者按（上下文推断，无外部依据）：由本节 (5.7) 的方程形状与 §17 的几何来源推得。]}
\begin{lemma}\label{st:5.3}
\textup{(i)} 取二元组 $(c,k)$，其中
\[
c\in\hom^0_{\A}(Y_1,Y_2),\qquad k:\hom_{\A}(Y_0,Y_1)\longrightarrow\hom_{\A}(Y_0,Y_2)[-1],
\]
满足
\begin{equation}\label{eq:5.7}
\begin{aligned}
\mu^1_{\A}(c)&=0,\\
\mu^1_{\A}(k(c_0))-k(\mu^1_{\A}(c_0))+\mu^2_{\A}(c,c_0)&=0\quad\text{对所有 }c_0\in\hom_{\A}(Y_0,Y_1).
\end{aligned}
\end{equation}
它确定模同态 $t:\mathcal T_{Y_0}(\Y_1)\to\Y_2$，满足
\[
[\mu^2_{\Q}(t,\iota)]=[\mathcal I^1(c)]\in\Hom_{H(\Q)}(\Y_1,\Y_2).
\]
\textup{(ii)} 给定 \textup{(i)} 中的二元组 $(c,k)$，以及任意
\[
\gamma\in\hom_{\A}(Y_1,Y_2)[-1],\qquad
\kappa:\hom_{\A}(Y_0,Y_1)\longrightarrow\hom_{\A}(Y_0,Y_2)[-2],
\]
令
\[
\begin{aligned}
\tilde c&=c+\mu^1_{\A}(\gamma),\\
\tilde k(c_0)&=k(c_0)+(-1)^{|c_0|-1}\bigl(\mu^1_{\A}(\kappa(c_0))-\kappa(\mu^1_{\A}(c_0))+\mu^2_{\A}(\gamma,c_0)\bigr),
\end{aligned}
\]
便得到\eref{5.7}的一个新解；与之相应的同态 $\tilde t$ 与原来的 $t$ 上同调。

\SourcePage{74}
\textup{(iii)} 设 $(c,k)$ 如 \textup{(i)}。此外，假设对每个 $X\in\Ob\A$，给定映射
\[
\begin{aligned}
d&:\hom_{\A}(X,Y_1)\longrightarrow\hom_{\A}(X,Y_2),\\
l&:\hom_{\A}(Y_0,Y_1)\otimes\hom_{\A}(X,Y_0)\longrightarrow\hom_{\A}(X,Y_2)[-1],\\
\xi&:\hom_{\A}(X,Y_1)\longrightarrow\hom_{\A}(X,Y_2)[-1],\\
\psi&:\hom_{\A}(Y_0,Y_1)\otimes\hom_{\A}(X,Y_0)\longrightarrow\hom_{\A}(X,Y_2)[-2],
\end{aligned}
\]
对所有 $c_0,b,c_1$ 满足
\begin{align*}
&\mu^1_{\A}(d(c_1))+d(\mu^1_{\A}(c_1))=0,\\
&d(\mu^2_{\A}(c_0,b))+\mu^1_{\A}(l(c_0,b))+(-1)^{|b|-1}l(\mu^1_{\A}(c_0),b)+l(c_0,\mu^1_{\A}(b))=0,\\
&\xi(\mu^1_{\A}(c_1))-\mu^1_{\A}(\xi(c_1))+d(c_1)-\mu^2_{\A}(c,c_1)=0,\\
&(-1)^{|b|-1}\psi(\mu^1_{\A}(c_0),b)+\psi(c_0,\mu^1_{\A}(b))-\mu^1_{\A}(\psi(c_0,b))\\
&\qquad+\xi(\mu^2_{\A}(c_0,b))+l(c_0,b)-\mu^2_{\A}(k(c_0),b)-\mu^3_{\A}(c,c_0,b)=0,
\end{align*}
并且下列链复形无上同调：
\begin{equation}\label{eq:5.8}
\begin{aligned}
D={}&(\hom_{\A}(Y_0,Y_1)\otimes\hom_{\A}(X,Y_0))[2]\\[-2pt]
&\quad\oplus\hom_{\A}(X,Y_1)[1]\oplus\hom_{\A}(X,Y_2),\\
\partial(c_0\otimes b,c_1,c_2)={}&\bigl((-1)^{|b|-1}\mu^1_{\A}(c_0)\otimes b+c_0\otimes\mu^1_{\A}(b),\\[-2pt]
&\quad\mu^1_{\A}(c_1)+\mu^2_{\A}(c_0,b),\ \mu^1_{\A}(c_2)+d(c_1)+l(c_0,b)\bigr).
\end{aligned}
\end{equation}
那么由 $(c,k)$ 定义的同态 $t$ 在 $H^0(\Q)$ 中诱导同构。事实上，只需对组成拟等价全子范畴 $\widetilde{\A}\subset\A$ 的一组对象验证这一条件，它就自动对其余所有对象成立。
\end{lemma}
\begin{proof}
(i) 直接计算表明，按下式定义 $t$，它便具有所有所需性质：
\[
\begin{aligned}
&t^d((c_0\otimes b,c_1),a_{d-1},\ldots,a_1)\\
&\quad=\mu^{d+1}_{\A}(c,c_1,a_{d-1},\ldots,a_1)+\mu^{d+1}_{\A}(k(c_0),b,a_{d-1},\ldots,a_1)\\
&\qquad+\mu^{d+2}_{\A}(c,c_0,b,a_{d-1},\ldots,a_1).
\end{aligned}
\]
(ii) 仿照 (i)，以 $(\gamma,\kappa)$ 代替 $(c,k)$，定义 $u\in\hom^{-1}_{\Q}(\mathcal T_{Y_0}\Y_1,\Y_2)$，便有 $\mu^1_{\Q}(u)=\tilde t-t$。

(iii) 要使 $[t]$ 为同构，必须使 $t^1$ 的映射锥无上同调。具体地，这些锥是链复形
\[
\begin{aligned}
\widetilde D={}&(\hom_{\A}(Y_0,Y_1)\otimes\hom_{\A}(X,Y_0))[2]\\[-2pt]
&\quad\oplus\hom_{\A}(X,Y_1)[1]\oplus\hom_{\A}(X,Y_2),\\
\tilde\partial(c_0\otimes b,c_1,c_2)={}&\bigl((-1)^{|b|-1}\mu^1_{\A}(c_0)\otimes b+c_0\otimes\mu^1_{\A}(b),\\[-2pt]
&\quad\mu^1_{\A}(c_1)+\mu^2_{\A}(c_0,b),\\[-2pt]
&\quad\mu^1_{\A}(c_2)+\mu^2_{\A}(c,c_1)+\mu^2_{\A}(k(c_0),b)+\mu^3_{\A}(c,c_0,b)\bigr).
\end{aligned}
\]
\SourcePage{75}
利用 $(\xi,\psi)$ 可以构造同构 $D\to\widetilde D$，即
\[
(c_0\otimes b,c_1,c_2)\longmapsto(c_0\otimes b,c_1,c_2+\xi(c_1)+\psi(c_0\otimes b)).
\]
所以\eref{5.8}无上同调蕴含所需结论。最后一项断言是显然的，因为限制函子 $\Q=\modcat(\A)\to\modcat(\widetilde{\A})$ 是拟等价（\stref{引理}{2.8}）。
\end{proof}

\begin{lemma}\label{st:5.4}
固定\stref{引理}{5.3}\textup{(i)} 中的 $(c,k)$。假设此外存在
\[
\begin{aligned}
g&=\sum_jg_2^j\otimes g_1^j\in(\hom_{\A}(Y_0,Y_1)\otimes\hom_{\A}(Y_2,Y_0))^1,\\
p&=\sum_jp_2^j\otimes p_1^j\in(\hom_{\A}(Y_0,Y_1)\otimes\hom_{\A}(Y_1,Y_0))^0,\\
q&=\sum_jq_2^j\otimes q_1^j\in\bigl(\hom_{\A}(Y_0,Y_1)\otimes\Hom_{\K}(\hom_{\A}(Y_0,Y_1),\hom_{\A}(Y_0,Y_0))\bigr)^{-1},
\end{aligned}
\]
满足
\[
\begin{aligned}
&\sum_j\bigl((-1)^{|g_1^j|-1}\mu^1_{\A}(g_2^j)\otimes g_1^j+g_2^j\otimes\mu^1_{\A}(g_1^j)\bigr)=0,\\
&\sum_jg_2^j\otimes\mu^2_{\A}(g_1^j,c)\\
&\quad=\sum_j(-1)^{|p_1^j|-1}\mu^1_{\A}(p_2^j)\otimes p_1^j+\sum_jp_2^j\otimes\mu^1_{\A}(p_1^j),\\
&\sum_jg_2^j\otimes\bigl(\mu^2_{\A}(g_1^j,k(\cdot))+\mu^3_{\A}(g_1^j,c,\cdot)\bigr)\\
&\quad=\sum_jq_2^j\otimes\mu^1_{\A}(q_1^j(\cdot))-\sum_jq_2^j\otimes q_1^j(\mu^1_{\A}(\cdot))\\
&\qquad+\sum_j(-1)^{|q_1^j|+|\cdot|-1}\mu^1_{\A}(q_2^j)\otimes q_1^j(\cdot)+\sum_jp_2^j\otimes\mu^2_{\A}(p_1^j,\cdot)\\
&\qquad+\sum_j b^j\otimes(e_{Y_0}\otimes b^{j,\vee}),
\end{aligned}
\]
其中，$e_{Y_0}$ 是 $Y_0$ 的恒等态射的代表，$b^j$ 为 $\hom_{\A}(Y_0,Y_1)$ 的一组基，$b^{j,\vee}$ 为其对偶基，而 $(e_{Y_0}\otimes b^{j,\vee})\in\Hom_{\K}(\hom_{\A}(Y_0,Y_1),\hom_{\A}(Y_0,Y_0))$ 为将 $b^j$ 映为 $e_{Y_0}$、将所有其他基元素映为零的线性映射。那么，$g$ 给出一次模同态 $u:\Y_2\to\hom_{\A}(Y_0,Y_1)\otimes\Y_0$，满足 $[\mu^2_{\Q}(u,t)]=[\pi]$，其中 $[\pi]$ 是\eref{5.3}中的映射。
\end{lemma}
\begin{proof}
定义预模同态
\[
\begin{aligned}
u&:\Y_2\longrightarrow\hom_{\A}(Y_0,Y_1)\otimes\Y_0,\\
u^d(b,a_{d-1},\ldots,a_1)&=\sum_jg_2^j\otimes\mu^{d+1}_{\A}(g_1^j,b,a_{d-1},\ldots,a_1),
\end{aligned}
\]
以及
\SourcePage{76}
\[
\begin{aligned}
v&:\mathcal T_{Y_0}(\Y_1)\longrightarrow\hom_{\A}(Y_0,Y_1)\otimes\Y_0,\\
v^d((c_0\otimes b,c_1),a_{d-1},\ldots,a_1)
&=\sum_jp_2^j\otimes\mu^{d+1}_{\A}(p_1^j,c_1,a_{d-1},\ldots,a_1)\\
&\quad+\sum_jq_2^j\otimes\mu^{d+1}_{\A}(q_1^j(c_0),b,a_{d-1},\ldots,a_1)\\
&\quad+\sum_jp_2^j\otimes\mu^{d+2}_{\A}(p_1^j,c_0,b,a_{d-1},\ldots,a_1)\\
&\quad-\sum_jg_2^j\otimes\mu^{d+2}_{\A}(g_1^j,c,c_1,a_{d-1},\ldots,a_1)\\
&\quad-\sum_jg_2^j\otimes\mu^{d+2}_{\A}(g_1^j,k(c_0),b,a_{d-1},\ldots,a_1)\\
&\quad-\sum_jg_2^j\otimes\mu^{d+3}_{\A}(g_1^j,c,c_0,b,a_{d-1},\ldots,a_1),
\end{aligned}
\]
两者次数分别为一与零。利用假设可算得 $\mu^1_{\Q}(u)=0$，并且 $w=\mu^2_{\Q}(u,t)-\mu^1_{\Q}(v)$ 为
\[
w^d((c_0\otimes b,c_1),a_{d-1},\ldots,a_1)=(-1)^{|c_0|}c_0\otimes\mu^{d+1}_{\A}(e_{Y_0},b,a_{d-1},\ldots,a_1).
\]
由与\stref{引理}{3.6}相同的论证，它与 $\pi$ 上同调。
\end{proof}

\begin{note}\label{st:5.5}
设处于\stref{引理}{5.3}的 (i)、(iii) 所述情形；另假设 $\A$ 对有限直和与平移封闭，且 $\Hom_{H(\A)}(Y_0,Y_1)$ 有限维，从而 $\Hom_{H(\A)}(Y_0,Y_1)\otimes Y_0$ 有意义。由构造，在 $H(\A)$ 中有下列正合三角：
\SourceEquationDiagram{diagram-5-9}{0.39046112}{5.9}
从这个角度看，扭转过程退居幕后：只需利用\stref{引理}{5.3}给出的一组条件，就能判定具有\eref{5.9}这种特殊形式的三角正合。不过要注意，三角中的第三个态射仅被间接刻画：它在 Yoneda 嵌入下的像是 $[\pi]\circ[t]^{-1}$。这时就需要\stref{引理}{5.4}。假设我们有该引理指定的额外数据，并考虑 $H(\A)$ 中的态射
\[
[g]=\sum_j[g_2^j]\otimes[g_1^j]:Y_2\longrightarrow\Hom_{H(\A)}(Y_0,Y_1)\otimes Y_0.
\]
在 Yoneda 嵌入下，它映为模同态 $[u]$，由构造满足 $[u]\circ[t]=[\pi]$。因此，\eref{5.9}中的第三个映射恰好就是 $[g]$。
\end{note}

\subsection{函子性}\label{sec:5d}
设 $\F:\A\to\B$ 为 $\Ai$ 范畴之间的上同调全忠实函子。取 $Y_0,Y_1\in\Ob\A$。
\begin{lemma}\label{st:5.6}
设两个扭转 $T_{Y_0}Y_1$ 与 $T_{\F(Y_0)}\F(Y_1)$ 都存在。那么在 $H^0(\B)$ 中有典范同构
\begin{equation}\label{eq:5.10}
\F(T_{Y_0}Y_1)\cong T_{\F(Y_0)}\F(Y_1).
\end{equation}
\end{lemma}
\SourcePage{77}
\begin{proof}
先给出一个较简单的初步论证。由\stref{推论}{3.16}，正合三角\eref{5.4}在 $H^0(\A)$ 中的同构意义下刻画 $T_{Y_0}(Y_1)$。由此结合\stref{推论}{3.9}，很容易推出\eref{5.10}中的两个对象同构。不过，这并未给出同构的典范选择；要得到更强的结论，必须稍微更仔细地论证。

沿用\eref{1.23}的思路，对任意 $Y\in\Ob\A$ 与 $\M\in\Ob(\modcat(\B))$，都有典范模同态
\[
s:\mathcal T_Y(\F^*\M)\longrightarrow\F^*(\mathcal T_{\F(Y)}\M).
\]
具体地，$s^1$ 将 $(c_0\otimes b,c_1)$ 映为 $((-1)^{|c_0|+|b|-1}c_0\otimes\F^1(b),(-1)^{|c_1|}c_1)$；对于 $d\geq2$，$s^d$ 将 $((c_0\otimes b,c_1),a_{d-1},\ldots,a_1)$ 映为 $((-1)^{\triangle}c_0\otimes\F^d(b,a_{d-1},\ldots,a_1),0)$，其中 $\triangle=|a_1|+\cdots+|a_{d-1}|+|b|+|c_0|-d$。考察 $d=1$ 就可见，$[s]$ 是 $H^0(\modcat(\A))$ 中的同构。\znote{69}{$s$ 比较的是“先扭转再沿 $\F$ 拉回”与“先沿 $\F$ 拉回再扭转”两个模，它在 $d=1$ 时已是 $H^0$ 中的同构，原因在于 $\F$ 上同调全忠实：$s^1$ 的两个分量分别由 $\F^1$ 与恒等给出，而 $\hom$ 空间上的同构正是全忠实性所保证的。高阶项 $s^d$ 的符号 $(-1)^{\triangle}$ 只是把 $\F^d$ 逐个移过前面的元素所生的 Koszul 符号，不影响同构性。[来源：并见 Seidel, Fukaya categories and Picard–Lefschetz theory, §5d（引理 5.6 的证明）。]}对 $Y=Y_0$、$\M=(\mathcal I_{\B}\circ\F)(Y_1)$ 应用这一点，得到
\begin{equation}\label{eq:5.11}
\mathcal T_{Y_0}(\F^*\circ\mathcal I_{\B}\circ\F)(Y_1)
\cong\F^*\bigl(\mathcal T_{\F(Y_0)}(\mathcal I_{\B}\circ\F)(Y_1)\bigr).
\end{equation}
由\eref{2.13}，\eref{5.11}的左边典范同构于 $\mathcal T_{Y_0}(\mathcal I_{\A}(Y_1))$；由 $T_{Y_0}$ 的定义，它典范同构于 $\mathcal I_{\A}(T_{Y_0}(Y_1))$，从而也典范同构于 $(\F^*\circ\mathcal I_{\B}\circ\F)(T_{Y_0}Y_1)$。对右边应用 $T_{\F(Y_0)}$ 的定义，可知它典范同构于 $(\F^*\circ\mathcal I_{\B})(T_{\F(Y_0)}\F(Y_1))$。所以，\eref{5.10}中的两个对象在施加 $\F^*\circ\mathcal I_{\B}$ 后成为典范同构的对象。再利用\eref{2.13}，可知 $\F^*\circ\mathcal I_{\B}$ 在 $\F$ 的像上上同调全忠实，因此在 $H^0(\B)$ 中同构于某个 $\F(X)$ 的对象所成的稍大对象类上，也上同调全忠实。\eref{5.10}的两边都属于这一类，故得到它们之间的典范同构。
\end{proof}
这样构造的同构关于 $Y_1$ 是自然的（事实上，再多作一些工作，就能将它延伸为函子范畴 $H^0(\operatorname{fun}(\A,\B))$ 中 $\F\circ T_{Y_0}$ 与 $T_{\F(Y_0)}\circ\F$ 之间的典范同构）。它也与典范映射 $Y_1\to T_{Y_0}Y_1$、$\F(Y_1)\to T_{\F(Y_0)}\F(Y_1)$ 相容。若进一步假设这些 $\Ai$ 范畴对有限直和与平移封闭，且有关的 $\Hom$ 空间有限维，那么正合三角\eref{5.4}在 $H(\F)$ 下的像，典范同构于对象 $\F(Y_0),\F(Y_1)$ 的相应三角。
